\section{Management de projet}

\subsection{Méthode de pilotage}

Le projet a été conduit selon une logique incrémentale inspirée du Kaizen :
avancer par versions courtes, tester chaque brique, intégrer progressivement,
puis documenter les décisions.

La priorité a été donnée à une démonstration fonctionnelle :

\begin{itemize}
    \item un flux utilisateur simple ;
    \item une capture audio réelle ;
    \item une reconnaissance musicale exploitable ;
    \item un affichage lisible ;
    \item une configuration réseau minimale.
\end{itemize}

\subsection{WBS -- Work Breakdown Structure}

\begin{lstlisting}
1. Gestion projet
   1.1 Cahier des charges
   1.2 Analyse fonctionnelle
   1.3 Planning du 11 mai au 11 juin
   1.4 Suivi des risques
   1.5 Documentation finale

2. Conception materielle
   2.1 Choix Raspberry Pi
   2.2 Choix e-paper
   2.3 Choix microphone
   2.4 Integration boutons
   2.5 Boitier / support style cassette
   2.6 Alimentation

3. Conception logicielle
   3.1 Architecture Python
   3.2 Capture audio
   3.3 Client AudD
   3.4 Controleur de detection
   3.5 Rendu e-paper
   3.6 Interface Flask
   3.7 Gestion Wi-Fi
   3.8 Nettoyage et documentation

4. Tests
   4.1 Test capture audio
   4.2 Test reconnaissance
   4.3 Test affichage
   4.4 Test boutons
   4.5 Test Wi-Fi
   4.6 Test demonstration complete

5. Rendus
   5.1 Dossier LaTeX
   5.2 Archives code
   5.3 Fichiers CAO et schemas
   5.4 Poster separe
   5.5 Soutenance
\end{lstlisting}

\subsection{Planning réalisé}

\begin{longtable}{|p{3cm}|p{5.2cm}|p{5.2cm}|}
\hline
\textbf{Période} & \textbf{Objectif} & \textbf{Résultat attendu} \\
\hline
11--15 mai & Cadrage du besoin et définition du concept. &
Cahier des charges initial, choix du concept S.Y.S.T.U.S. \\
\hline
16--20 mai & Architecture et choix techniques. &
Découpage Python, choix Raspberry Pi, écran e-paper, Flask et reconnaissance AudD. \\
\hline
21--25 mai & Capture audio et reconnaissance. &
Module \texttt{AudioRecorder}, génération WAV, client \texttt{AudDClient}. \\
\hline
26--30 mai & Affichage et états système. &
Écrans \textit{idle}, \textit{listening}, \textit{thinking}, \textit{result}, QR code. \\
\hline
31 mai--4 juin & Wi-Fi et interaction physique. &
Interface Flask, scan réseau, configuration \texttt{nmcli}, boutons GPIO. \\
\hline
5--9 juin & Intégration et stabilisation. &
Boucle principale, modes \texttt{setup}/\texttt{running}, nettoyage des états. \\
\hline
10 juin & Finalisation documentaire. &
Dossier LaTeX, archives projet, préparation dépôt FabManager. \\
\hline
11 juin & Soutenance. &
Support oral, démonstration prototype, questions-réponses. \\
\hline
\end{longtable}

\subsection{Kanban de fin de projet}

\begin{longtable}{|p{3cm}|p{7cm}|p{3cm}|}
\hline
\textbf{État} & \textbf{Tâches} & \textbf{Priorité} \\
\hline
Terminé & Structure Python du projet. & Haute \\
\hline
Terminé & Capture audio avec \texttt{arecord}. & Haute \\
\hline
Terminé & Client AudD pour reconnaissance de fichier WAV. & Haute \\
\hline
Terminé & États de détection et rendu e-paper. & Haute \\
\hline
Terminé & Interface web Flask de configuration Wi-Fi. & Haute \\
\hline
Terminé & Gestion de deux boutons GPIO : mode Wi-Fi et lancement d'écoute. & Haute \\
\hline
Terminé & Modélisation 3D du boîtier, du support et de la pièce écran. & Haute \\
\hline
À tester sur cible & Démonstration complète sur Raspberry Pi Zero 2 W monté. & Haute \\
\hline
À compléter & Prix exacts et nomenclature finale. & Moyenne \\
\hline
À compléter & Photos, schémas définitifs et fichiers CAO. & Moyenne \\
\hline
Hors délai & Optimisation autonomie et boîtier final industrialisable. & Basse \\
\hline
\end{longtable}

\subsection{Matrice des risques}

\begin{longtable}{|p{3.2cm}|p{2cm}|p{2cm}|p{3.2cm}|p{3cm}|}
\hline
\textbf{Risque} & \textbf{Prob.} & \textbf{Impact} & \textbf{Prévention} & \textbf{Plan de secours} \\
\hline
Micro mal configuré & Moyenne & Élevé & Tester tôt la commande \texttt{arecord}. & Utiliser un fichier WAV de test pour la démo logicielle. \\
\hline
API indisponible & Moyenne & Élevé & Prévoir une connexion stable et des exemples connus. & Montrer les états système et un résultat pré-enregistré. \\
\hline
Wi-Fi instable & Moyenne & Moyen & Interface de configuration et test \texttt{nmcli}. & Connexion manuelle avant soutenance. \\
\hline
E-paper lent & Haute & Moyen & Limiter les rafraîchissements aux changements d'état. & Démontrer uniquement les états essentiels. \\
\hline
Boutons avec rebond & Moyenne & Moyen & Ajouter un verrou et une temporisation logicielle. & Lancer l'action depuis le code pendant la démo. \\
\hline
Secret API exposé & Moyenne & Moyen & Déplacer le token vers une variable d'environnement. & Révoquer le token après rendu et documenter la limite. \\
\hline
Temps de finition court & Haute & Moyen & Prioriser le flux complet avant les finitions. & Présenter les limites clairement dans le bilan. \\
\hline
\end{longtable}

\subsection{Journal de décisions}

\begin{longtable}{|p{2.5cm}|p{4cm}|p{4cm}|p{3cm}|}
\hline
\textbf{Date} & \textbf{Décision} & \textbf{Justification} & \textbf{Impact} \\
\hline
Mai 2026 & Utiliser Python. & Rapidité de prototypage et nombreuses bibliothèques. & Développement plus rapide. \\
\hline
Mai 2026 & Utiliser Flask pour le Wi-Fi. & Léger, simple, suffisant pour une interface locale. & Interface web minimale opérationnelle. \\
\hline
Mai 2026 & Utiliser AudD pour la reconnaissance. & API simple à intégrer dans le délai. & Dépendance au réseau et à un token. \\
\hline
Mai 2026 & Utiliser l'e-paper Waveshare 4.2 V2. & Affichage sobre, lisible et adapté au style objet. & Rafraîchissement lent mais acceptable. \\
\hline
Juin 2026 & Séparer setup et running. & Clarifie l'expérience utilisateur. & Deux modes simples à expliquer en soutenance. \\
\hline
Juin 2026 & Déclencher l'écoute par clic bouton. & Évite une écoute continue et donne le contrôle à l'utilisateur. & Une nouvelle reconnaissance nécessite un nouvel appui. \\
\hline
\end{longtable}
